\section{Ecuaciones Matemáticas}

\LaTeX{} con \texttt{amsmath} ofrece el sistema más avanzado mundial para matemáticas académicas \cite{amsmath_manual}.

\subsection{Ecuación Simple}

El bloque \texttt{equation} genera ecuaciones numeradas automáticamente:

\begin{equation}
E = mc^2 \label{eq:einstein_real}
\end{equation}

Para generar la ecuación \eqref{eq:einstein_real}, se utiliza el siguiente código:

\begin{verbatim}
\begin{equation}
E = mc^2 \label{eq:einstein}
\end{equation}
\end{verbatim}


\subsection{Ecuaciones Múltiples Alineadas}

Es posible alinear varias ecuaciones en un solo bloque usando \texttt{align}. Las ecuaciones se alinean alrededor del símbolo \texttt{\&} y se numeran automáticamente \cite{amsmath_manual}:

\begin{align}
\nabla \cdot \mathbf{E} &= \frac{\rho}{\epsilon_0} \label{eq:gauss_real} \\
\nabla \cdot \mathbf{B} &= 0 \label{eq:gauss_mag_real} \\
\nabla \times \mathbf{E} &= -\frac{\partial \mathbf{B}}{\partial t} \label{eq:faraday_real} \\
\nabla \times \mathbf{B} &= \mu_0\mathbf{J}+\mu_0\epsilon_0\frac{\partial\mathbf{E}}{\partial t} \label{eq:ampere_real}
\end{align}

El ejemplo anterior muestra las ecuaciones de Maxwell (Eqs.~\eqref{eq:gauss_real}--\eqref{eq:ampere_real}) que se pueden escribir con el siguiente código \cite{amsmath_manual}:

\begin{verbatim}
\begin{align}
\nabla \cdot \mathbf{E} &= 
\frac{\rho}{\epsilon_0} \label{eq:gauss} \\
\nabla \cdot \mathbf{B} &= 
0 \label{eq:gauss_mag} \\
\nabla \times \mathbf{E} &= 
-\frac{\partial \mathbf{B}}{\partial t} 
\label{eq:faraday} \\
\nabla \times \mathbf{B} &= 
\mu_0\mathbf{J}+\mu_0\epsilon_0
\frac{\partial\mathbf{E}}
{\partial t} \label{eq:ampere}
\end{align}
\end{verbatim}

Como se puede observar, existen comandos específicos para símbolos matemáticos avanzados como \texttt{\textbackslash nabla} (del operador nabla), \texttt{\textbackslash cdot} (producto escalar), \texttt{\textbackslash times} (producto vectorial), y operadores de derivada parcial \texttt{\textbackslash partial}, y fracciones \texttt{\textbackslash frac} \cite{amsmath_manual}.

También es posible suprimir la numeración de líneas intermedias usando \texttt{\textbackslash notag} o \texttt{\textbackslash nonumber}:

\begin{align}
P &= \frac{1}{2\pi} \int_0^{2\pi} f(\theta) \, d\theta \notag \\
&= \lim_{n\to\infty} \frac{1}{n} \sum_{k=1}^n f\left(\frac{2\pi k}{n}\right) \label{eq:promedio_real} \\
&= \mathbb{E}[f(X)] \notag
\end{align}

El código para generar la ecuación anterior \eqref{eq:promedio_real} es el siguiente:

\begin{verbatim}
\begin{align}
P &= \frac{1}{2\pi} \int_0^{2\pi} 
f(\theta) \, d\theta \notag \\
&= \lim_{n\to\infty} \frac{1}{n} \sum_{k=1}^n 
f\left(\frac{2\pi k}{n}\right) \label{eq:promedio} \\
&= \mathbb{E}[f(X)] \notag
\end{align}
\end{verbatim}


\subsection{Matrices y Sistemas Lineales}

\LaTeX también permite escribir matrices de forma sencilla usando entornos como \texttt{bmatrix} para matrices con corchetes \cite{amsmath_manual}:

\begin{equation}
\mathbf{A}\mathbf{x} = \mathbf{b}, \quad
\mathbf{A} = \begin{bmatrix}
1 & 2 & 0 \\
3 & 0 & 1 \\
0 & 2 & 4
\end{bmatrix},
\quad \mathbf{x} = \begin{bmatrix} x_1 \\ x_2 \\ x_3 \end{bmatrix}
\label{eq:sistema_matrices_real}
\end{equation}

La ecuación anterior \eqref{eq:sistema_matrices_real} se puede escribir con el siguiente código:

\begin{verbatim}
\begin{equation}
\mathbf{A}\mathbf{x} = \mathbf{b}, \quad
\mathbf{A} = \begin{bmatrix}
1 & 2 & 0 \\
3 & 0 & 1 \\
0 & 2 & 4
\end{bmatrix},
\quad \mathbf{x} = \begin{bmatrix} 
x_1 \\ x_2 \\ x_3 \end{bmatrix}
\label{eq:sistema_matrices}
\end{equation}
\end{verbatim}

\subsection{Ecuaciones con Casos y Subecuaciones}

Utilizando el entorno \texttt{cases}, es posible escribir funciones definidas por casos \cite{amsmath_manual}:

\begin{equation}
f(x) = \begin{cases}
x^2 & x \geq 0 \\
-x^2 & x < 0
\end{cases} \label{eq:casos_real}
\end{equation}

El siguiente código genera la función definida por casos de la ecuación \eqref{eq:casos_real}:

\begin{verbatim}
\begin{equation}
f(x) = \begin{cases}
x^2 & x \geq 0 \\
-x^2 & x < 0
\end{cases} \label{eq:casos}
\end{equation}
\end{verbatim}

Por último, el entorno \texttt{subequations} permite agrupar varias ecuaciones bajo una numeración común \cite{amsmath_manual}:

\begin{subequations}\label{eq:sistema_completo_real}
\begin{align}
\frac{dx}{dt} &= ax + by \label{eq:x_real} \\
\frac{dy}{dt} &= cx + dy \label{eq:y_real}
\end{align}
\end{subequations}

El sistema de ecuaciones (\ref{eq:x_real})--(\ref{eq:y_real}) se puede escribir con el siguiente código:

\begin{verbatim}
\begin{subequations}\label{eq:sistema_completo}
\begin{align}
\frac{dx}{dt} &= ax + by \label{eq:x} \\
\frac{dy}{dt} &= cx + dy \label{eq:y}
\end{align}
\end{subequations}
\end{verbatim}



